% Glossary terms from ch21: lines of the form \item[Term] Definition.
\item[Model predictive control (MPC)] A feedback controller that, at every step, solves a finite-horizon optimal control problem from the measured state, applies the first input and repeats; constraints are part of the optimization rather than an afterthought.
\item[Receding horizon] The moving window of $N$ future steps over which MPC optimizes; after each step the window is shifted by one step and the plan is recomputed.
\item[Condensed QP] The form of the MPC problem in which the predicted states are eliminated with $\mat{X}=\mat{S}_x\state_0+\mat{S}_u\mat{U}$, leaving a dense quadratic program in the stacked inputs $\mat{U}$ alone.
\item[Prediction matrices] $\mat{S}_x$ and $\mat{S}_u$, built from powers of $\mat{A}$ and $\mat{B}$, that map the current state and the stacked inputs to the stacked predicted states.
\item[Quadratic program (QP)] Minimization of a convex quadratic function $\tfrac12\vect{z}\T\mat{H}\vect{z}+\vect{f}\T\vect{z}$ subject to linear inequality constraints; the problem MPC with a linear model, quadratic cost and linearized constraints solves at every step.
\item[Active constraint] An inequality constraint that holds with equality at the solution and has a positive Lagrange multiplier; it is the constraint that shapes the solution.
\item[Linearized avoidance constraint] The half-plane $\vect{n}_k\T(\pos_k-\hat{\vect{o}}_k)\ge r_k$ that replaces the non-convex disk constraint $\norm{\pos_k-\hat{\vect{o}}_k}\ge r_k$ at step $k$, with the normal $\vect{n}_k$ taken from the previous plan; it is convex and conservative.
\item[Soft constraint] A constraint $g(\vect{z})\le s$ with a slack variable $s\ge0$ that is penalized in the cost; it keeps the problem feasible and gives the least violation when the hard version cannot be satisfied.
\item[Exact penalty] A linear slack weight $w$ larger than every Lagrange multiplier of the hard problem, for which the soft solution coincides with the hard one whenever the latter is feasible.
\item[Terminal cost] The weight $\state_N\T\mat{P}\state_N$ on the last predicted state, usually the LQR cost-to-go from the Riccati equation, which stands in for the cost beyond the horizon and underlies the stability proof of MPC.
\item[Recursive feasibility] The property that feasibility of the MPC problem at one step implies feasibility at the next; with a terminal set and a static world it follows from the shifted plan.
\item[Warm start] Initializing the QP solver with the previous plan shifted by one step, which typically halves the number of iterations.
\item[ADMM] The alternating direction method of multipliers, a first-order splitting method for convex problems; the OSQP iteration used by the book's QP solver.
\item[Chance constraint] A constraint required to hold with probability at least $1-\delta$; for a Gaussian predicted intruder it becomes the half-plane constraint with the radius inflated by $\kappa_\delta\sqrt{\vect{n}_k\T\mat{\Sigma}_k\vect{n}_k}$.
\item[Formation and communication-range constraints] Requirements $\norm{\pos^i_k-\pos^j_k-\vect{d}_{ij}}\le e_{\max}$ and $\norm{\pos^i_k-\pos^j_k}\le R_{\mathrm{comm}}$ on pairs of drones, convex disks around a moving center that MPC enforces through inscribed polygons.
